\documentclass[12pt]{article}
\usepackage[utf8]{inputenc}
\usepackage[T1]{fontenc}
\usepackage{hyperref}
\title{Pwaamei Language Resources for K-12 Curriculum Development}
\author{Compiled by Claude Sonnet 5.5 (claude-sonnet-5-5) for Andr\'as Kornai}
\date{}
\begin{document}
\maketitle

\section{Introduction}
Pwaamei (also spelled Poamei; ISO 639-3 pme) is an Austronesian Kanak language of New Caledonia, spoken in the commune of Voh, North Province (Hoot Ma Waap customary area). Speaker figures in the sources I found vary from about 219 to 290 (2009 census), down from 325 reported in 1974. It is listed as definitely endangered, children are beginning not to learn it, and one source notes no literacy in the language.

\section{Primary Educational Resources}
No primers, textbooks or school programmes for Pwaamei were found in my searches.

\section{Linguistic and Reference Materials}
\begin{itemize}
\item Wikipedia overview: \url{https://en.wikipedia.org/wiki/Pwaam%C3%A8i_language}
\item Overview of New Caledonian languages: \url{https://en.wikipedia.org/wiki/New_Caledonian_languages}
\item Endangered Languages Project entry: \url{https://endangeredlanguages.com/lang/2526?hl=en}
\item Verbix endangered-languages entry: \url{https://docs.verbix.com/EndangeredLanguages/Pwaamei}
\end{itemize}

\section{Religious and Community Resources}
Joshua Project reports that no translated or oral Scripture in Pwaamei is available: \url{https://joshuaproject.net/people_groups/14418}.

\section{Digital and Audio Resources}
None found.

\section{Curriculum Development Context}
Little documentation was found. Any curriculum would require fieldwork and community partnership with Voh speakers, and possibly links to New Caledonia's Kanak language teaching framework, which I did not research.
\end{document}